\documentclass{jsarticle}
\usepackage[dvipdfmx, final]{graphicx}
\usepackage{amsmath}
\usepackage{amssymb}
\usepackage{chemfig}
\begin{document}
\title{}
\author{}
\date{}
\maketitle

   稼働機材の燃料エネルギー、消費量、
   政府の計画、企業の予測、IBMとMicrosoftとの予測の取り方、
   


   トランプ大統領の政府の予測、計画、Microsoftも間を考えている。

   建物の中、公共施設の外の、椅子に、手軽に使える、機材、

   50年先の企業の目標、
   1500年、5000年先の予測(ここだけ)
   不老長寿も、視野に入れている
   苫米地博士の予見力の本の趣旨を題目にしての、インタビュー、
   それも、大から中、小に至る、話、
   Quntam computer on toplogy, lang straight on view,
   ここは、博士は200年先と言っていた。

##ここは、関係ない 
## トポロジは、サーストン・ペレルマン多様体や一般相対性理論の場の理論の
## ヒルベルト空間になっている。
## 低次元多様体が、京都大学の先生が特に研究している。結び目理論や
## Jones多項式には、高次元多様体から、低次元多様体へと、
## 広中先生の数値分解にもなっている。
## アインシュタイン博士の上を行く研究が、結構幾多にもある。
##
##   トポロジは、位相幾何学の一分野であり、多様体や基本群、
##   射像や同型、準同型、幾何学の不変量を見つけたり、
##

   Copilot on windows machine,
   ここは、既に実現している、
   publicate in chair and build in, hologram display,
   ここは、既に、アップルがしている。もっと、上でしている。
   発想が逆転している製品を、Vision proでして、


   私らも、休憩しながら、飲み物でも飲んで、先を行くとしよう。
\end{document}


